\section{Limitations}

The study is limited to Vietnamese text, the fixed operational label set, and the evaluated domains and splits. The idiom/figurative-language alpha is only 0.5746, indicating substantial annotation ambiguity for that phenomenon. Macro-F1 can conceal class-specific errors, particularly for scarce positive labels. Three-run variation remains limited. For Q1a, the nominal paired-bootstrap 95\% CI for the macro-pragmatic gap relative to the XLM-R-large pragmatic-only baseline is [+0.07,+1.20] pp, whereas all six head-level intervals cross zero. We therefore avoid head-specific or family-wise significance claims; these results do not substitute for a larger statistical study. Using the same Q1a training settings reduces optimization confounding, but the systems differ in several auxiliary training components; the comparison therefore does not isolate the causal effect of any single component. Q3 is non-monotonic and does not establish a general sample-efficiency claim. Calibration depends on the specified split, bins, and probability source. External UIT-VSFC, UIT-VSMEC, and AIVIVN evaluations test affective transfer of the polarity/emotion outputs under their respective protocols, not complete six-label pragmatic transfer. Finally, the training-only rationale decoder is not a faithful explanation module at inference.
